% theme: scale_drawings_and_angle_bisectors
\MajorQuestionLesson{縮図と角の二等分線の定理}
\SetRowSpace{3mm}


\Needspace{80mm}
\TypeQuestion{}{similarity_application_map_scale}
\FigProblemBlock{27mm}{fig/lesson_06_map_scale}{
\Text{縮尺 $1:25000$ の地図上で、地点 $\mathrm{A}$ から地点 $\mathrm{B}$ を通って地点 $\mathrm{C}$ まで進む道を考える。地図上の $\Lseg{AB}$ の長さは $3.6\mathrm{cm}$、$\Lseg{BC}$ の長さは $2.4\mathrm{cm}$ である。ただし、図は略図である。}
\QQ{実際の $\Lseg{AB}$ の道のりを $\mathrm{m}$ で求めなさい。}{}
}{
\QQ{地点 $\mathrm{A}$ から地点 $\mathrm{B}$ を通って地点 $\mathrm{C}$ まで進む実際の道のりを $\mathrm{km}$ で求めなさい。}{}
\QQ{同じ道を縮尺 $1:10000$ の地図に表すと、$\Lseg{AB}$ と $\Lseg{BC}$ を合わせた地図上の長さは何 $\mathrm{cm}$ になるか求めなさい。}{}
}

\Needspace{80mm}
\TypeQuestion{}{similarity_application_river_scale_drawing}
\FigProblemBlock{53mm}{fig/lesson_06_river_scale}{
\Text{川岸の2点 $\mathrm{A},\mathrm{B}$ と、対岸の点 $\mathrm{P}$ の位置を調べ、縮尺 $1:400$ の縮図をかいた。点 $\mathrm{P}$ から直線 $\Lseg{AB}$ に引いた垂線の足を $\mathrm{H}$ とする。縮図上の $\Lseg{AB}$ の長さは $4.5\mathrm{cm}$、$\Lseg{PH}$ の長さは $2.8\mathrm{cm}$ である。川の両岸は平行で、川幅は一定とする。}
\QQ{実際の川幅 $\Lseg{PH}$ を $\mathrm{m}$ で求めなさい。}{}
}{
\QQ{実際の $\Lseg{AB}$ の長さを $\mathrm{m}$ で求めなさい。}{}
\QQ{同じ場所を縮尺 $1:1000$ の縮図に表すと、線分 $\Lseg{PH}$ の長さは何 $\mathrm{cm}$ になるか求めなさい。}{}
}
